% Source: 02 - Tue Sep 29 2026.pdf, page 3. Content remains in source order.
\sourcepage{02}{3}
An encoding turns an abstract decision problem $\Pi_A$ into a concrete decision problem:
\[
\begin{array}{rcl}
\Pi_A:\mathcal I&\longrightarrow&\{0,1\}\\
\downarrow e(\cdot)&&\\
\Pi_C:\bits^*&\longrightarrow&\{0,1\}.
\end{array}
\]
The two problems are related as follows:
\[
\forall x\in\bits^*:\quad
\text{(1)}\quad\Pi_C(x)=1\Longleftrightarrow
\exists i\in\mathcal I:(x=e(i))\land(\Pi_A(i)=\mathrm{yes}).
\]

\textbf{NOTE 1:} $\Pi_C(x)$ is defined on all strings in $\bits^*$. From $(1,\Rightarrow)$ we have:
\[
\forall x\in\bits^*:\quad\nexists i\in\mathcal I:x=e(i)\ \Rightarrow\ \Pi_C(x)=0.
\]
\textbf{NOTE 2:} The size of a concrete instance is $|x|$
